\documentclass[11pt,titlepage]{report}
\usepackage{Preamble}

\newtheorem{theorem}{Theorem}
\newtheorem{corollary}[theorem]{Corollary}
\newtheorem{lemma}[theorem]{Lemma}
\newtheorem{claim}[theorem]{Claim}
\newtheorem{definition}[theorem]{Definition}
\newtheorem{fact}[theorem]{Fact}
\newenvironment{proof}{{\bf Proof\quad}}{\hfill\rule{2mm}{2mm}}
\DeclareMathOperator*{\cond}{cond_2}

\begin{document}

\begin{titlepage}
    \newgeometry{margin=3cm}
	\centering
    \includegraphics[width=0.5\linewidth]{images/EPFL.png}\\[0.25cm] 	% University Logo
    \textsc{\LARGE École Polytechnique Fédérale de Lausanne}\\ \vspace{\fill}
    \textbf{\textsc{\fontsize{30}{30}\selectfont Super-Resolution Off the Grid}}\\ \vspace{\fill}		
	\textsc{\LARGE EPFL - Semester Project in Computer Science}\\[0.4cm]
	\rule{\linewidth}{0.2 mm} \\[0.5 cm]
	Sébastien Ollquist \\[2cm] \today
\end{titlepage}
\restoregeometry

\thispagestyle{numberonly}
\begin{summary}
\section*{Abstract}
Super-Resolution is the tool that allows us to upgrade the quality of images. The goal of this project is to study a recent paper on super-resolution and discuss the main algorithm in order to eventually find an improvement using information theoretical bounds.
\end{summary}

\begin{fquote}[Albert Einstein]
    In theory, theory and practice are the same. In practice, they are not.
\end{fquote}

\chapter{Introduction}
With the increasing amount of data nowadays, it has become extremely important to develop tools that allow us to treat it better. In the case of images, unfortunately a lot of them are still taken with technology that still misses some important information. Recently, a new theory called super-resolution has emerged. Its aim is to, given access to some coarse measurements, recover estimates of some point sources located in a $d$-dimensional plane.\par 

The paper we will study here is called "Super-Resolution Off the Grid" and has been published by Qingqing Huang and Sham M. Kakade in September 2015. From a theoretical point of view, super-resolution is the problem of recovering a superposition of point sources using (noise corrupted) bandlimited measurements. In essence, suppose we are given a collection of points that lie on the $d$-dimensional plane, to which we have no access whatsoever. We use measurements from the Fourier domain and super-resolve the point sources so that the estimates we get are the best possible ones.\par 

The goal of this project is to get familiar with the theory of super-resolution along with all important mathematical preliminaries and study the main result presented in the above mentioned paper. The idea is then to compare this work with other related ones and discuss the efficiency of the main algorithm. Ultimately, it would be a nice addition to start working on an improvement of the main algorithm in order to reduce the sample complexity from quadratic space to the information theoretical linear one.
\chapter{Mathematical refresher}
\section{A mathematical theory of super-resolution}
\paragraph{Mathematical abstraction} We consider $k$ point sources in $d$ dimensions, where the points are separated by a distance at least $\Delta$ (in Euclidean distance). The $d$-dimensional signal $x(t)$ can be modeled as a weighted sum of $k$ Dirac measures in $\mathbb{R}^d$ as $$x(t)=\sum_{j=1}^k w_j\delta_{\mu^{(j)}},$$ where the $\mu^{(j)}$'s are the point sources in $\mathbb{R}^d$ and $w_j\in\mathbb{C}$ the weights such that $|w_j|<C$ for every $j\in[k]$ and some absolute constant $C>0$. $\Delta$ is by definition the minimal possible distance between any two pair of points of the $d$-dimensional plane. Formally, we regard $\Delta$ in terms of Euclidean distance as $$\Delta=\min_{j\neq j'}||\mu^{(j)}-\mu^{(j')}||_2.$$
\paragraph{Measurement function} There are two main factors that represent our signal: the weights $w_j$ and the complex low pass point spread functions $e^{i\pi\langle s,t\rangle}$ for some measurement $s$. The former is simply a weighting factor that is not too relevant to us here. The latter however represents the response of an imaging system to a point source. That is, it determines the overall performance level of our system and will be of crucial interest to our study.\par 
We define the measurement function $f(s):\mathbb{R}^d\to\mathbb{C}$ as being the convolution of the point source $x(t)$ with the point spread function $e^{i\pi\langle s,t\rangle}$. Formally, $$f(s)=\sum_{j\in[k]}w_je^{i\pi\langle\mu^{(j)},s\rangle}.$$ The function $f$ is in essence a sum of weighted Fourier coefficients. Note that the measurements $s$ in the noisy setting are corrupted with a uniformly bounded perturbation $z$ so that the noisy recovery problem becomes $$\widetilde{f}(s)=f(s)+z(s),$$ in which $|z(s)|\leq\epsilon_z$ for every $s$ and a constant $\epsilon_z\in(0,1/2)$. The idea is that given access to the signal $x(t)$, we wish to generate a set of random bandlimited Fourier measurements and evaluate the noisy function $\widetilde{f}$ on every measurement $s$ in order to recover the parameters $\{w_j,\mu^{(j)}:j\in[k]\}$ of the point source signal as best as we can. The problem statement will be described in more detail in Chapter 4.

\section{Main tools}
The main mathematical tools that are needed to understand the paper essentially lie around the subject of linear algebra. They include notions related to vectors, matrices and tensors. We also require some probabilistic analysis tools since the algorithm is partly random. In this chapter, we introduce those tools, prove the more important results and closely relate them to the paper.
\subsection{Generalised eigenvalue problem}
Before introducing the generalised eigenvalue problem, it is important to recall what a condition number is for a particular matrix. Suppose for instance we have a matrix $X\in\mathbb{R}^{m\times n}$. We let $\lambda_1,\ldots,\lambda_n$ be the eigenvalues of $X^TX$ (with repetitions) and arrange them so that $\lambda_1\geq\ldots\geq\lambda_n\geq 0$. Then, the $\sigma_1\geq\ldots\geq\sigma_n\geq 0$ such that $\sigma_i=\sqrt{\lambda_i}$ are called the \textit{singular values} of $X$. Define $\sigma_{max}(X)=\sigma_1$ and $\sigma_{min}(X)=\sigma_n$. We then define the condition number of a matrix to be the ratio between the largest and the smallest singular value of $X$. That is 
\begin{equation}
    \cond(X)=\sigma_1/\sigma_n=\frac{\sigma_{max}(X)}{\sigma_{min}(X)}.
\end{equation}
The above factor governs the noise tolerance of the generalised eigenvalue problem, i.e. it is the measure of sensitiveness of a matrix to arbitrary perturbations. Taking its limiting value will allow us to state whether or not the main algorithm achieves stable recovery of the point sources.\par
The goal of an eigenvalue problem is to simply find the eigenvalues of a particular matrix $A$. To do so, we generally solve the equation $AU=UV$, where $U$ is an orthogonal matrix. Multiplying each side by $U^T$ yields $$AUU^T=UVU^T\Rightarrow A=UVU^T=UVU^{-1},$$ where $\mathbb{R}^{d\times d}\ni U:=(u_1,\ldots,u_d)$ are the eigenvectors and $\mathbb{R}^{d\times d}\ni V:=Diag\left[(\lambda_1,\ldots,\lambda_d)^T\right]$ the eigenvalues. However, in the generalised version of this problem, we add another random matrix $B$ such that the problem becomes 
\begin{flalign*}
    AU=BUV&\Rightarrow AUU^T=BUVU^T\\
    &\Rightarrow A=BUVU^T=BUVU^{-1},
\end{flalign*}
where $U$ again are the eigenvectors and $V$ the eigenvalues. Note that $A$ and $B$ are both symmetric and that $B$ is generated randomly with the only condition being that it is positive definite. We form the pair $(A,B)$ called the \textit{pencil} and the pair $(U,V)$ called the \textit{eigenpair}. Therefore, we introduce a first version of what is called the matrix pencil method, in which given a pair of matrices $(A,B)$, we wish to find the generalised eigenvalues $\lambda$ for which there is a vector ${\bf x}$ such that $A{\bf x}=\lambda B{\bf x}$. Notice that the eigenvalues of $A$ are the solution to the generalised eigenvalue problem where $B=I$.
\subsection{Common norms}
\paragraph{Vectors} Norms are a powerful linear algebra tool that allows us to retreive important information regarding a particular matrix or vector. The most common vector norm is the $L^p$-norm defined for a vector $X$ of size $n$ as 
\begin{equation}
    ||X||_p=\left(\sum_{i\in[n]}x_i^p\right)^{1/p}.
\end{equation}
Note that $p$ can take any value in $\mathbb{N}^*$ but the most common ones are 1, 2 or $\infty$. In essence, the $L^1$-norm is the sum of absolute values of our vector, the $L^2$-norm represents the Euclidean distance of the vector from the origin and the $L^\infty$ norm is by definition the largest element in the vector, in absolute value.
\paragraph{Matrices} The same way we define a norm operator for vectors, we can define a similar operator for matrices, that will however have a slightly different meaning. Essentially, a matrix norm is a vector norm in a vector space whose elements are matrices. Let $A$ be a matrix of size $m\times n$. We define the following important norms for $A$:
\begin{flalign}
    ||A||_1&=\max_{j\in[n]}\sum_{i\in[m]}|A_{i,j}|\\
    ||A||_2&=\sqrt{\lambda_{max}(A^*A)}=\sigma_{max}(A)\label{eqn:spectral_norm}\\
    ||A||_\infty&=\max_{i\in[m]}\sum_{j\in[n]}|A_{i,j}|,
\end{flalign} where $\lambda_{max}(A^*A)$ is the largest eigenvalue of $A^*A$ in which $A^*$ denotes the conjugate transpose of $A$. Note that equality \ref{eqn:spectral_norm} is sometimes called the \textit{spectral norm} of a matrix. We also consider another important matrix norm called the \textit{Frobenius} norm, which is an equivalent version of the vector $L^2$-norm. It essentially represents the size of the matrix $A$, and we define it as 
\begin{equation}
    ||A||_F=\sqrt{\sum_{i\in[m]}\sum_{j\in[n]}A_{i,j}^2}.
\end{equation}
\subsection{Gershgorin's disk theorem}
A powerful result of linear algebra allows us to relate the eigenvalues of some matrix $A$ to points in the complex plane. Namely, the Gershgorin's disk theorem gives us a pretty good and efficient way to find the eigenvalues of a particular matrix by drawing disks in the complex plane and relating them to the eigenvalues. For the purpose, let $Y\in\mathbb{C}^{k\times k}$ be a complex valued matrix. For every $j\in[k]$, let $$\mathcal{D}_j=\big\{x\in\mathbb{C}:||x-Y_{j,j}||\leq\sum_{j'\neq j}|Y_{j,j'}|\big\}$$ be the disk associated with the $j$'th row of $Y$, defined on the complex plane. Then, Gershgorin's disk theorem states that all the eigenvalues of $Y$ must lie within the union of all the above defined disks. The intuition behind this theorem starts with the definition of an eigenvalue. Suppose for instance that $\lambda$ is an eigenvalue of $Y$. Then we know we can write $Yv=\lambda v$, where $v$ is an eigenvector of $Y$. Select $j$ to be the index whose entry has largest magnitude in $v$. Writing the matrix-vector product as a sum yields $[Yv]_j=\sum_{j'\in[n]}Y_{j,j'}v_{j'}=\lambda v_j$. We can split this sum into two distinct sums as $\sum_{j'\in[n]}Y_{j,j'}v_{j'}=Y_{j,j}v_j+\sum_{j'\neq j}Y_{j,j'}v_{j'}$. Putting the factors of $v_{j'}$ on the same side gives us $$\sum_{j'\neq j}Y_{j,j'}v_{j'}=(\lambda-Y_{j,j})v_j.$$ We can divide by $v_j$ on both sides and take the absolute value to obtain $$\bigg|\sum_{j'\neq j}Y_{j,j'}v_{j'}/v_j\bigg|=\big|\lambda-Y_{j,j}\big|.$$ Using the triangle inequality for absolute values, we get
\begin{flalign*}
    \big|\lambda-Y_{j,j}\big|&\leq\sum_{j'\neq j}\big|Y_{j,j'}v_{j'}/v_j\big|\\
    &\leq\sum_{j'\neq j}\big|Y_{j,j'}\big|\big|v_{j'}/v_j\big|\text{ since $|ab|=|a||b|$ for any $a,b$}\\
    &\leq\sum_{j'\neq j}\big|Y_{j,j'}\big|,
\end{flalign*} where in the last line we used the fact that $v_j$ has the largest magnitude in $v$ such that $v_j\geq v_{j'}$ for every $j$. Hence, the Gershgorin disk centered at some diagonal $j$ has radius always at most the sum of all non-diagonal entries of $Y$ of the same row $j$.
\subsection{Probabilistic inequalities}
Here we present the main probability bounds that are used to prove the correctness of the main algorithm. The first lemma is a generalisation of the Hoeffding bound to random matrices. It tells us that with high probability, the values of a random Hermitian matrix concentrate around the mean. The second considers the projection of a high dimensional random vector onto a lower dimension plane. Namely, suppose that we consider a random vector that lies on the $m$-dimensional unit sphere and another fixed vector. Then, the lemma states that with high probability, the absolute value of their dot product is large. We now state the aforementionned lemmas:
\begin{lemma} (Matrix Hoeffding)
    Let $\{X^{(1)},\ldots,X^{(m)}\}$ be a set of independent, random, Hermitian matrices of dimension $k\times k$, with identical distribution $X$. Assume that $E[X]$ is finite, and $X^2\preceq\sigma^2 I$ for some positive constant $\sigma$ almost surely. Then, for all $\epsilon\geq 0$, $$Pr\left(\left|\left|\frac{1}{m}\sum_{i=1}^m X^{(i)} - E[X]\right|\right|_2\geq\epsilon\right)\leq ke^{-\frac{m^2\epsilon^2}{8\sigma^2}}.$$
\end{lemma}
\begin{lemma} (Vector random projection)
    Let $a\in\mathbb{R}^m$ be a random vector distributed uniformly over the $m$-dimensional unit sphere and fix $v\in\mathbb{C}^m$. Then, for $\delta\in(0,1)$, $$Pr\left(|\langle a,v\rangle|\leq\frac{||v||_2}{\sqrt{em}}\delta\right)\leq\delta.$$
\end{lemma}
\begin{proof} (for vector random projection)
    We give an intuitive proof of the above lemma, which is mostly based on lemma 2.2 from \cite{proofJLLemma}. Let $X_1,\ldots,X_d$ be independent samples from a standard Gaussian distribution and let $Y=(X_1,\ldots,X_d)/||X||$ be a vector formed by the normalised Gaussians. It is straightforward to see that $Y$ is uniformly distributed on the $d$-dimensional sphere. Let now $Z\in\mathbb{R}^k$ be the projection of $Y$ onto the first $k$ coordinates and let $L=||Z||^2$. Observe that $E[L]=d/k$ for $d<k$. Then, it is true that 
    \begin{align}
        Pr\big(L\leq\beta k/d\big)\leq\beta e^{\frac{k}{2}(1-\beta)},\quad\beta<1\\
        Pr\big(L\geq\beta k/d\big)\leq\beta e^{\frac{k}{2}(1-\beta)},\quad\beta>1
    \end{align}
    Both bounds together tell us that the squared spectral norm of the projection is in fact fairly concentrated around its expectation. Now to apply the above to our lemma, set $L$ to be $|\langle a,v\rangle|^2$ so that the bound becomes $$Pr\left(|\langle a,v\rangle|\leq\frac{\sqrt{\beta}||v||_2}{\sqrt{m}}\right)\leq\beta e^{\frac{k}{2}(1-\beta)}.$$
    Setting $\beta=\delta^2/e$ yields 
    $$Pr\left(|\langle a,v\rangle|\leq\frac{||v||_2}{\sqrt{em}}\delta\right)\leq\frac{\delta^2}{e}e^{\frac{k}{2}(1-\delta^2/e)}.$$
    Hence to prove the bound, it suffices to show that 
    $\frac{\delta^2}{e}e^{\frac{k}{2}(1-\delta^2/e)}\leq\delta$ or equivalently $\delta e^{\frac{k}{2}(1-\delta^2/e)}\leq e.$
    Since by assumption $\delta<1$ and $1<e^{k(1+\delta^2/e)/2}$, the proof is completed.
\end{proof}

\section{Tensor decomposition}
A tensor is a generalisation of a matrix to more than two dimensions. We can think of a tensor as a point in $\mathbb{C}^{m_1\times\ldots\times m_p}$ where $p$ is the order of the tensor. Most of the time here, $p=3$ since three dimensions suffice for our analysis. Note that if $T$ is an order three tensor of dimensions $m\times n\times p$, we can view it as a collection of $p$ matrices of size $m\times n$ stacked on top of each other. For example, the entry $A_{i,j,k}$ of a 3-tensor $A$ will simply be the $(i,j)$'th entry of the $k$'th matrix \cite{tensorMethods}. Furthermore, we define the \textit{rank} of a tensor $T$ as the minimum $k$ such that we can write $T$ as the sum of rank one tensors.
\subsection{Canonical polyadic decomposition}
Canonical Polyadic (CP) decomposition (or \textit{tensor rank decomposition}) is a procedure for tensors equivalent to that of the SVD for matrices in which we wish to express a $p$-tensor as an outer product of $p$ vectors. We first recall that the outer product of two vectors $u\in\mathbb{C}^m$ and $v\in\mathbb{C}^n$ gives a matrix $M\in\mathbb{C}^{m\times n}$ such that $[M]_{r,c}=u_rv_c$ for $r\in[m]$ and $c\in[n]$. From there, suppose that we are working with a rank $k$ 3-tensor $V\in\mathbb{C}^{m\times n\times p}$. CP decomposition allows us to express $V$ as
\begin{equation}
    V=\sum_{n\in[k]}[A]_{:,n}\circ[B]_{:,n}\circ[C]_{:,n},    
\end{equation} where $\circ$ denotes the outer product, $A\in\mathbb{C}^{m\times k}$, $B\in\mathbb{C}^{n\times k}$, $C\in\mathbb{C}^{p\times k}$ and where $[M]_{:,n}$ denotes the vector formed by the $n$'th column of $M$. We then say that the tensor $V$ admits the decomposition $V=A\otimes B\otimes C$, where $\otimes$ denotes the tensor product operation. Equivalently, element-wise we have 
\begin{equation}
    V_{j_1,j_2,j_3}=\sum_{n\in[k]}A_{j_1,n}B_{j_2,n}C_{j_3,n},
\end{equation} for $j_1\in[m],j_2\in[n],j_3\in[p]$.
\subsection{Tensor decomposition as a multilinear mapping}
We can use an alternative definition of tensor decomposition using the notion of a multi-linear mapping. Namely, for given dimensions $m_A$, $m_B$, $m_C$, the mapping $V(\cdot,\cdot,\cdot):\mathbb{C}^{m\times m_A}\times\mathbb{C}^{m\times m_B}\times\mathbb{C}^{m\times m_C}\to\mathbb{C}^{m_A\times m_B\times m_C}$ is defined as: 
\begin{equation}
    \left[V(X_A,X_B,X_C)\right]_{i_1,i_2,i_3}=\sum_{j_1,j_2,j_3\in[m]}V_{j_1,j_2,j_3}[X_A]_{j_1,i_1}[X_B]_{j_2,i_2}[X_C]_{j_3,i_3}.
\end{equation}
We can verify that for a particular vector $a\in\mathbb{C}^m$, the projection $V(I,I,a)$ of $V$ along the 3rd dimension is $V(I,I,a)=ADiag(C^T a)B^T$ as long as $V$ admits a tensor decomposition $V=A\otimes B\otimes C$. Indeed,
\begin{flalign*}
    \left[V(I,I,a)\right]_{i_1,i_2}&=\sum_{j_1,j_2,j_3\in[m]}V_{j_1,j_2,j_3}[I]_{j_1,i_1}[I]_{j_2,i_2}[a]_{j_3}\\
    &=\sum_{j_1,j_2,j_3\in[m]}\sum_{n\in[k]} A_{j_1,n}B_{j_2,n}C_{j_3,n}[I]_{j_1,i_1}[I]_{j_2,i_2}[a]_{j_3}\\
    &=\sum_{j_3\in[m]}\sum_{n\in[k]} A_{i_1,n}B_{i_2,n}C_{j_3,n}[a]_{j_3}\\
    &=\underbrace{\sum_{n\in[k]}A_{i_1,n}B_{i_2,n}}_{=AB}\underbrace{\sum_{j_3\in[m]}C_{j_3,n}[a]_{j_3}}_{=Diag(C^T a)}\\
    &=\left[ADiag(C^T a)B^T\right]_{i_1,i_2}.
\end{flalign*} Since the above is true for every $i_1,i_2$, and we do not consider the $i_3$ coordinate since $a$ is a vector, we get the desired equality.\par 
In a more general case, we can consider the projection of the tensor $V$ on three $m$ column matrices $X$, $Y$ and $Z$. Again we suppose that $V$ admits the decomposition $A\otimes B\otimes C$. In that case, it is true that 
\begin{flalign*}
    [V(X,Y,Z)]_{i_1,i_2,i_3}&=\sum_{j_1,j_2,j_3\in[m]}V_{j_1,j_2,j_3}[X]_{j_1,i_1}[Y]_{j_2,i_2}[Z]_{j_3,i_3}\\
    &=\sum_{j_1,j_2,j_3\in[m]}\sum_{n\in[k]}A_{j_1,n}B_{j_2,n}C_{j_3,n}[X]_{j_1,i_1}[Y]_{j_2,i_2}[Z]_{j_3,i_3}\\
    &=\sum_{j_1,j_2,j_3\in[m]}\sum_{n\in[k]}[X^T]_{i_1,j_1}A_{j_1,n}[Y^T]_{i_2,j_2}B_{j_2,n}[Z^T]_{i_3,j_3}C_{j_3,n}\\
    &=\sum_{n\in[k]}[X^TA]_{i_1,n}[Y^TB]_{i_2,n}[Z^TC]_{i_3,n}\\
    &=[(X^TA)\otimes(Y^TB)\otimes(Z^TC)]_{i_1,i_2,i_3}.
\end{flalign*}
Hence, we conclude that $V(X,Y,Z)$ admits the decomposition $(X^TA)\otimes(Y^TB)\otimes(Z^TC)$ for any arbitrary matrices $X$, $Y$ and $Z$.
\subsection{Tensor Frobenius norm}
We can also extend the matrix norm definition to that of a tensor. However in our work, we will only need the equivalent Frobenius norm which will be useful in the main algorithm. For the purpose, let $V$ be a rank $k$ $p$-tensor of dimensions $m\times n\times p$. Then, the Frobenius norm of $V$ is simply the root of the sum of all squares, i.e.
\begin{equation}
    ||V||_F=\sqrt{\sum_{i\in[m]}\sum_{j\in[n]}\sum_{l\in[p]}[V]_{i,j,l}^2}.
\end{equation}
\subsection{Jennrich's algorithm for tensor decomposition}
As previously stated, the goal of tensor decomposition is to, given a tensor $T$ of rank $k$, decompose it as a sum of rank 1 tensors of appropriate dimensions. Jennrich's algorithm is commonly used for tensor decomposition. We denote the following theorem:
\begin{theorem}
    Let $T=\sum_{j=1}^k u_i\otimes v_i\otimes w_i$ be a tensor in which each set of vectors $\{u_i\}_i$ and $\{v_i\}_i$ are linearly independent. Moreover, each pair of vectors in $\{w_i\}_i$ are also linearly independent. Then, the above decomposition is unique up to rescaling, and there is an efficient algorithm to find it.
\end{theorem}
The aforementioned algorithm is described in Algorithm \ref{alg:Jennrich}.
\begin{algorithm}
    \caption{Jennrich's algorithm for tensor decomposition}
    \label{alg:Jennrich}
    \begin{algorithmic}
        \State {\bf Input}: a tensor $\widetilde{F}\in\mathbb{C}^{m\times m\times 3}$ of rank $k$.
        \State Choose random unit vectors $a,b\in\mathbb{R}^m$.
        \State Compute $\widetilde{F}(I,I,a)=UD_aV^T$, where $D_a=Diag(\langle w^{(i)},a\rangle)$.
        \State Compute $\widetilde{F}(I,I,b)=UD_bV^T$, where $D_b=Diag(\langle w^{(i)},b\rangle)$.
        \State Compute the diagonalisations $\widetilde{F}(I,I,a)\widetilde{F}(I,I,b)^{-1}$ and $\widetilde{F}(I,I,b)\widetilde{F}(I,I,a)^{-1}$.
        \State Solve the linear system to recover the $w_j$'s.
        \State {\bf Return} $U$, $V$, $W$.
    \end{algorithmic}
\end{algorithm}\par
Observe that $$\widetilde{F}(I,I,a)\widetilde{F}(I,I,b)^{-1}=UD_aV^T(V^T)^{-1}D_b^{-1}U^{-1}=UD_aD_b^{-1}U^{-1},$$ and similarly $$\widetilde{F}(I,I,a)^{-1}\widetilde{F}(I,I,b)=(V^T)^{-1}D_a^{-1}U^{-1}UD_bV^T=(V^T)^{-1}D_a^{-1}D_bV^T.$$ The correctness of the algorithm follows from the uniqueness of eigendecomposition of a matrix when the eigenvalues are distinct. For a random choice of $a$ and $b$ (in our case we choose the basis vectors $e_1$ and $e_2$), with high probability the eigenvalues are unique so we can recover the $u_i$'s and the $v_i$'s easily by simply recovering the columns of $U$ and $V$, respectively \cite{algoToolbox}.
\chapter{Initial work}
\section{Prony's method}
In 1795, the french mathematician Gaspard de Prony came up with a very useful trick to solve a recovery problem, which aims to reconstruct functions from their values at given points. Essentially, we can view those functions as an $n$'th degree polynomial from which we want to recover the roots. This is today called \textit{Prony's method} and can easily be applied to the theory of super-resolution \cite{pronyMethod}.
\paragraph{General form} In a general sense, consider functions of the form $z=f(x)=\sum_{j\in[n]}\mu_j\rho_j^x$, where $\mu_j\in\mathbb{R}$ and $\rho_j\in\mathbb{R}_+$ for every $1\leq i,j\leq n$. We want to recover those parameters from the samples $z_k:=f(x_k)$ of $f$. We are faced with what is called Prony's condition, which states that given an undetermined number of measurements or sampling points $z_0,z_1,\ldots$, we impose that $$(z*p)_k=\sum_{j=0}^n z_{k+1}p_j=0,$$ for $k=0,1,2,\ldots$ and we wish to solve the above equations for the $p_j$'s. We see that the solutions are given by $p(x)=(x-\widehat{\rho}_1)\cdots(x-\widehat{\rho}_n)$. In a more algorithmic sense, we seek to 
\begin{enumerate}
    \item solve the above linear system for $p$ and 
    \item find the zeroes of $p(x)$ to obtain the $\rho_j$'s.
\end{enumerate}
\paragraph{Link with super-resolution} Prony's method can in fact be applied to many domains. For example in microscopy, images often contain brightness spikes on the plane but the optical system that contains all the hardware to create the picture essentially acts as a low pass filter: it will keep the low frequencies of a signal and attenuate the frequencies that are higher than a pre-determined cutoff frequency. By doing so, those spikes are turned into what are called \textit{point spread functions}, that determine the performance of an imaging system. Similarly to our problem, the goal of Prony's method is to reconstruct the image by recovering estimates of the original spikes.\par 
We can view the plane of an image as part of $\mathbb{R}^2$. Let its spikes be at positions $X\subset\mathbb{R}^2$ and let them have amplitude $a_x$. Note that $X$ is sparse. Taking the Fourier transform yields $$z_\alpha:=z_{\alpha_1,\alpha_2}=\sum_{x\in X}a_xe^{\alpha_1 x_1+\alpha_2 x_2},$$ where $\alpha=(\alpha_1,\alpha_2)\in\mathbb{Z}^2$. Through the low pass filter, we obtain a finite amount of measurements $$z_\alpha=\sum_{x\in X}a_x\rho_x^\alpha,$$ for $|\alpha_1|,|\alpha_2|\leq n$ where $\rho_x^\alpha=\rho_{x,1}^{\alpha_1}\rho_{x,2}^{\alpha_2}$ and $\rho_{x,j}=e^{ix_j}$ for $j=1,2$. Note that the recovery of the $z_\alpha$'s is exactly the 2D version of Prony's problem: the $a$'s are our weights and the $\rho_x^\alpha$'s are our point sources.
\section{Univariate case: Matrix-pencil method}
As an extension of Prony's method which is not stable, we consider the main algorithm in the univariate case, the setting of which is as follows. We construct two $m\times m$ complex valued Hankel matrices $H_0$ and $H_1$, that is, matrices such that their skew-diagonals\footnote{A skew-diagonal is the diagonal in the North-East direction.} are constants. We have $D_w\in\mathbb{C}^{k\times k}$ where $[D_w]_{j,j}=w_j$ and $D_\mu\in\mathbb{C}^{k\times k}$ where $[D_\mu]_{j,j}=e^{i\pi\mu^{(j)}}$. We furthermore have a Vandermonde matrix $V_m\in\mathbb{C}^{m\times k}$ defined as $$V_m=
\begin{pmatrix}
    1&\ldots&1\\
    \big(e^{i\pi\mu^{(1)}}\big)^1&\ldots&\big(e^{i\pi\mu^{(k)}}\big)^1\\
    \vdots&\ldots&\vdots\\
    \big(e^{i\pi\mu^{(1)}}\big)^{m-1}&\ldots&\big(e^{i\pi\mu^{(k)}}\big)^{m-1}
\end{pmatrix}$$ We construct a 3rd order tensor $F\in\mathbb{C}^{m\times m\times 2}$ in which $F_{i,i',j}=[H_{j-1}]_{i,i'}$ for $j=1,2$ and $i,i'\in[m]$. The following computations will help us better understand tensor decomposition and the main algorithm. We first denote the following fact:
\begin{fact}
    In the setting above, $F$ admits the unique rank $k$ tensor decomposition $F=V_m\otimes V_m\otimes (V_2D_w)$.
\end{fact}
\begin{proof}
    We start by computing the product $V_2D_w$. By definition of $V_m$, we have that $[V_2]_{r,c}=\big(e^{i\pi\mu^{(c)}}\big)^{r-1}$ for $r\in[2]$ and $c\in[k]$. Multiplying with the diagonal matrix $D_w$ yields $$[V_2D_w]_{r,c}=w_c\big(e^{i\pi\mu^{(c)}}\big)^{r-1},\quad\forall r\in[2],c\in[k].$$
    We aim to show that 
    \begin{flalign*}
        F_{i_1,i_2,i_3}&=\sum_{n=1}^k[V_m]_{i_1,n}[V_m]_{i_2,n}[V_2D_w]_{i_3,n}\\
        &=\sum_{n=1}^k\big(e^{i\pi\mu^{(n)}}\big)^{i_1-1}\big(e^{i\pi\mu^{(n)}}\big)^{i_2-1}w_n\big(e^{i\pi\mu^{(n)}}\big)^{i_3-1}\\
        &=\sum_{n=1}^k w_n\big(e^{i\pi\mu^{(n)}}\big)^{i_1+i_2+i_3-3},
    \end{flalign*} for $i_1,i_2\in[m]$ and $i_3\in[2]$. It can be verified that the measurements that form both Hankel matrices are given by $$f(s)=\sum_{j\in[k]}w_j\big(e^{i\pi\mu^{(j)}}\big)^s,$$ with $s=0,\ldots,2m-1$ for $H_0$ and $s=1,\ldots,2m$ for $H_1$. For the details of the proof, see Appendix \ref{appendix:hankel}. Then, since $F_{i_1,i_2,i_3}=[H_{i_3-1}]_{i_1,i_2}$, for a fixed $i_3$ we have either $F_{i_1,i_2,1}=[H_0]_{i_1,i_2}$ or $F_{i_1,i_2,2}=[H_1]_{i_1,i_2}$. Note that since we are dealing with Hankel matrices, $[H_0]_{i_1,i_2}=f(i_1+i_2-2)$. Hence,
    $$F_{i_1,i_2,i_3}=\sum_{n\in[k]}w_n\big(e^{i\pi\mu^{(n)}}\big)^{i_1+i_2+i_3-3}=f(i_1+i_2+i_3-3),$$ as required.
\end{proof}\par 
Recall that Jennrich's algorithm computes the simultaneous diagonalisations of $F(I,I,v_1)$ and $F(I,I,v_2)$, in which $v_1$ and $v_2$ are two random projections of $\mathbb{R}^m$. Note that setting $v_1=e_1$ and $v_2=e_2$ yields $H_0$ and $H_1$ exactly. Hence, we conclude that the matrix-pencil method studied here is just a special case of Jennrich's algorithm where both projections are done on the two basis vectors $e_1$ and $e_2$.
\section{Multivariate toy case}
This section describes a preview of the main algorithm when the measurements are not random and where we assume for simplicity that the point sources are uniformly distributed in $[-1,+1]$. We follow the definitions given in the paper. We have $D_w=Diag(w_j)\in\mathbb{C}^{k\times k}$ and $[V_d]_{r,c}=e^{i\pi\mu_r^{(c)}}\in\mathbb{C}^{d\times k}$ the latter in which $\mu_n^{(i)}$ are i.i.d. $\sim Unif([-1,+1])$. Furthermore, we have $F_{n_1,n_2,n_3}=f(s)\big|_{s=e_{n_1}+e_{n_2}+e_{n_3}},$ for all $n_1,n_2,n_3\in[d]$. We have the following fact:
\begin{fact}
    In the setting above, $F$ admits the tensor decomposition $F=V_d\otimes V_d\otimes(V_dD_w)$.
\end{fact}
\begin{proof}
    We wish to show that $f(s)\big|_{s=s^{(n_1)}+s^{(n_2)}+s^{(n_1)}}=\sum_{j=1}^k w_je^{i\pi(\mu_{n_1}^{(j)}+\mu_{n_2}^{(j)}+\mu_{n_3}^{(j)})}$. To do so, we first compute the matrix product $V_dD_w$. Since $[V_d]_{r,c}=e^{i\pi\mu_r^{(c)}}$ for $r\in[d],c\in[k]$ and $D_w=Diag(w_j)$, we directly have that $$[V_dD_w]_{r,c}=w_ce^{i\pi\mu_r^{(c)}},\quad r\in[d],c\in[k],$$ so that by definition of tensor decomposition, we get
    \begin{flalign*}
        F_{n_1,n_2,n_3}&=\sum_{j=1}^k[V_d]_{n_1,j}[V_d]_{n_2,j}[V_dD_w]_{n_3,j}\\
        &=\sum_{j=1}^k e^{i\pi\mu_{n_1}^{(j)}}e^{i\pi\mu_{n_2}^{(j)}}w_je^{i\pi\mu_{n_3}^{(j)}}\\
        &=\sum_{j=1}^k w_je^{i\pi(\mu_{n_1}^{(j)}+\mu_{n_2}^{(j)}+\mu_{n_3}^{(j)})}\\
        &=f(s)\big|_{s=s^{(n_1)}+s^{(n_2)}+s^{(n_1)}}=f(e_{n_1}+e_{n_2}+e_{n_3}),
    \end{flalign*} as required.
\end{proof}\par
The key idea here is to notice that the entries $e^{i\pi\mu_n^{(j)}}$ are i.i.d. and uniformly distributed over the unit circle of the complex plane. This is due to the fact that the $\mu_n^{(j)}$ factors are uniform. This implies that the factor $V_d$ almost surely has full column rank and therefore that the tensor decomposition exists and is unique (since each vector contained in the tensor's matrices are linearly independent). Again by uniformity assumption, the condition number of $V_d$ concentrates around 1 so the above algorithm achieves stable recovery. We will see in more details in the next chapter.
\chapter{Main algorithm}
Let us recall what the goal of the problem is. Broadly speaking, super-resolution aims to recover a superposition of point sources using bandlimited measurements that may be corrupted with noise. In this setting, the created algorithm works in the Fourier domain where each of the $k$ points of the $d$-dimensional plane are separated by a distance at least $\Delta$. Hence, the frequencies of the Fourier measurements are bounded by $O(1/\Delta)$. The idea of the procedure is to take random bandlimited measurements with cutoff frequency bounded by $\Omega(\sqrt{d}/\Delta)$ and perform Tensor decomposition to recover the estimates of the point sources.\par 
This chapter discusses the main algorithm along with its analysis. We will go through all the important results and describe in detail their implication on the main procedure to effectively recover the point sources.
\section{Algorithm description}
\paragraph{Input and output} The algorithm takes as input a cutoff frequency $R$, the defined number of measurements $m$ and a noisy measurement function $\tilde{f}(\cdot)$ and outputs the set of estimates $$\{\widehat{w}_j,\widehat{\mu}^{(j)}:j\in[k]\},$$ where the $\widehat{w}_j$'s are the complex weight coefficients and the $\widehat{\mu}^{(j)}$'s are the estimates of the point sources. Note that in case the noise is non-existant (i.e. when $\epsilon_z=0$), the parameters are recovered exactly. Otherwise, stable recovery implies that the estimates are such that $$\min_\pi\max\left\{||\widehat{\mu}^{(j)}-\mu^{(\pi(j))}||_2:j\in[k]\right\}\leq poly(d,k)\epsilon_z.$$ In words, the estimates of the point sources differ from their real values by at most the noise $\epsilon_z$ scaled by a polynomial function that depends on the number of dimensions $d$ and the number of point sources $k$.
\paragraph{Measurements} We first generate a set $\mathcal{S}=\{s^{(1)},\ldots,s^{(m+n+1)}\}$ where $s^{(1)},\ldots,s^{(m)}$ are $m$ i.i.d. samples from the Gaussian distribution $\mathcal{N}(0,R^2I_{d\times d})$, $s^{(m+n)}=e_n$ for all $n\in[d]$ and $s^{(m+n+1)}=0$ and we let $m'=m+d+1$. We then take a sample $v$ from the unit sphere and set $v^{(1)}=v$ and $v^{(2)}=2v$. Since we are in dimension $d$, the sample $v$ is such that $x_1^2+\ldots+x_d^2=1$. Finally, we construct a tensor $\tilde{F}\in\mathbb{C}^{m'\times m'\times 3}:\tilde{F}_{n_1,n_2,n_3}=\tilde{f}(s)|_{s=s^{(n_1)}+s^{(n_2)}+v^{(n_3)}}$.
\paragraph{Tensor decomposition} We apply Jennrich's algorithm on $\tilde{F}$ to obtain the estimates $\widehat{V}_{S'}$ and $\widehat{D}_w$. Note that the estimate $\widehat{V}_{S'}$ is normalised so that its last line is all ones, in order to match the definition of $V_{S'}$ (see next section).
\paragraph{Read of estimates} We finish by recovering the real part of the estimates of the point sources by setting $\widehat{\mu}^{(j)}=Real(\log([\widehat{V}_S]_{[m+1:m+d,j]})/(i\pi))$ for every $j\in[k]$ and we find the best possible matching coefficients by setting $\widehat{W}=\arg\min_{W\in\mathbb{C}^k}||\widehat{F}-\widehat{V}_{S'}\otimes\widehat{V}_{S'}\otimes \widehat{V}_dD_w||_F$, where $||\cdot||_F$ is the Frobenius norm of $(\cdot)$.
\section{Stability guarantee}
The main guarantee on the stability of the recovery of the point sources is given by $$\min_\pi\max\left\{||\widehat{\mu}^{(j)}-\mu^{(\pi(j))}||_2:j\in[k]\right\}\leq C\frac{\sqrt{d}k^5}{\Delta\delta_v}\frac{w_{max}}{w_{min}^2}\left(\frac{1+2\epsilon_x}{1-2\epsilon_x}\right)^{5/2}\epsilon_z,$$ where $C$ is a universal constant such that the bound holds with probability at least $(1-\delta_s)$ over the random sampling of $\mathcal{S}$ and with probability at least $(1-\delta_v)$ over the random projections in Jennrich's algorithm. Note that {\bf the spatial resolution of an imaging device may be measured by how closely lines can be resolved}.
\section{Analysis}
% First of all, note that the sample complexity is determined solely by the condition numbers of $F$ and $V_S$.
Let us first recall what variables are used in the main algorithm.
\begin{itemize}
    \item $[V_S]_{r,c}=e^{i\pi\langle\mu^{(c)},s^{(r)}\rangle}$ is the \textit{characteristic matrix}.
    \item $[V_d]_{r,c}=e^{i\pi\mu_r^{(c)}}$ is the matrix that contains all the complex valued coefficients that lie on the unit circle of the complex plane due to the uniformity of the $\mu$'s.
    \item $V_{S'}=\big[V_S,\, V_d,\, [1]^k\big]^T$ is the matrix that is part of the decomposition we are trying to compute.
    \item $$V_2=
    \begin{pmatrix}
        e^{i\pi\langle\mu^{(1)},v^{(1)}\rangle}&\ldots&e^{i\pi\langle\mu^{(k)},v^{(1)}\rangle}\\
        e^{i\pi\langle\mu^{(1)},v^{(2)}\rangle}&\ldots&e^{i\pi\langle\mu^{(k)},v^{(2)}\rangle}\\
        1&\ldots&1
    \end{pmatrix}.$$
\end{itemize}
\begin{fact}
    In the exact case (i.e. when $\epsilon_z=0$), the constructed tensor $F$ admits a rank-$k$ decomposition $$F=V_{S'}\otimes V_{S'}\otimes (V_2D_w).$$
\end{fact}
\begin{proof}
    As usual, we start by computing the product $V_2D_w$. Since $D_w$ is a diagonal matrix, we have $$V_2D_w=
    \begin{pmatrix}
        w_1e^{i\pi\langle\mu^{(1)},v^{(1)}\rangle}&\ldots&w_ke^{i\pi\langle\mu^{(k)},v^{(1)}\rangle}\\
        w_1e^{i\pi\langle\mu^{(1)},v^{(2)}\rangle}&\ldots&w_ke^{i\pi\langle\mu^{(k)},v^{(2)}\rangle}\\
        w_1&\ldots&w_k
    \end{pmatrix}.$$
    For simplicity of notation, define $v^{(3)}=0$ so that $e^{i\pi\langle\mu^{(j)},v^{(3)}\rangle}=1$ for every $j\in[k]$. Using this, we can express the $r$'th line and $c$'th column in $V_2D_w$ as $$[V_2D_w]_{r,c}=w_ce^{i\pi\langle\mu^{(c)},v^{(r)}\rangle},$$ for every $r\in[3]$ and $c\in[k]$. What we want to show is $$F_{n_1,n_2,n_3}=\sum_{n\in[k]}[V_{S'}]_{n_1,n}[V_{S'}]_{n_2,n}[V_2D_w]_{n_3,n},$$ where $F_{n_1,n_2,n_3}=\tilde{f}(s)\big|_{s=s^{(n_1)}+s^{(n_2)}+v^{(n_3)}}$ are the noisy measurements.
\end{proof}
\subsection{Exact case recovery}


\chapter{Discussion}

\chapter{Improvement}

\chapter{Conclusion}

\appendix
\chapter{Construction of Hankel matrices}
\label{appendix:hankel}
In this section, we prove that the Hankel matrices $H_0$ and $H_1$ defined earlier admit the respective diagonalisations $V_mD_wV_m^T$ and $V_mD_wD_\mu V_m^T$. To do so, we first compute the product $D_wV_m^T$. We have 
\begin{flalign*}
    D_wV_m^T&=
    \begin{pmatrix}
        w_1&0&\ldots&0\\
        0&w_2&\ldots&0\\
        \vdots&\vdots&\ddots&\vdots\\
        0&0&\ldots&w_k
    \end{pmatrix}
    \begin{pmatrix}
        1&e^{i\pi\mu^{(1)}}&\ldots&\big(e^{i\pi\mu^{(1)}}\big)^{m-1}\\
        1&e^{i\pi\mu^{(2)}}&\ldots&\big(e^{i\pi\mu^{(2)}}\big)^{m-1}\\
        \vdots&\vdots&\ddots&\vdots\\
        1&e^{i\pi\mu^{(k)}}&\ldots&\big(e^{i\pi\mu^{(k)}}\big)^{m-1}
    \end{pmatrix}\\
    &=\begin{pmatrix}
        w_1&w_1e^{i\pi\mu^{(1)}}&\ldots&w_1\big(e^{i\pi\mu^{(1)}}\big)^{m-1}\\
        w_2&w_2e^{i\pi\mu^{(2)}}&\ldots&w_2\big(e^{i\pi\mu^{(2)}}\big)^{m-1}\\
        \vdots&\vdots&\ddots&\vdots\\
        w_k&w_ke^{i\pi\mu^{(k)}}&\ldots&w_k\big(e^{i\pi\mu^{(k)}}\big)^{m-1}
    \end{pmatrix},
\end{flalign*} so that
\begin{flalign*}
    H_0=V_mD_wV_m^T&=
    \begin{pmatrix}
        1&\ldots&1\\
        e^{i\pi\mu^{(1)}}&\ldots&e^{i\pi\mu^{(k)}}\\
        \vdots&\cdots&\vdots\\
        \big(e^{i\pi\mu^{(1)}}\big)^{m-1}&\ldots&\big(e^{i\pi\mu^{(k)}}\big)^{m-1}
    \end{pmatrix}
    \begin{pmatrix}
        w_1&w_1e^{i\pi\mu^{(1)}}&\ldots&w_1\big(e^{i\pi\mu^{(1)}}\big)^{m-1}\\
        w_2&w_2e^{i\pi\mu^{(2)}}&\ldots&w_2\big(e^{i\pi\mu^{(2)}}\big)^{m-1}\\
        \vdots&\vdots&\ddots&\vdots\\
        w_k&w_ke^{i\pi\mu^{(k)}}&\ldots&w_k\big(e^{i\pi\mu^{(k)}}\big)^{m-1}
    \end{pmatrix}\\
    &=\begin{pmatrix}
        \sum_{j\in[k]}w_j&\sum_{j\in[k]}w_je^{i\pi\mu^{(1)}}&\ldots&\sum_{j\in[k]}w_j\big(e^{i\pi\mu^{(1)}}\big)^{m-1}\\
        \sum_{j\in[k]}w_je^{i\pi\mu^{(1)}}&\sum_{j\in[k]}w_j\big(e^{i\pi\mu^{(1)}}\big)^2&\ldots&\sum_{j\in[k]}w_j\big(e^{i\pi\mu^{(1)}}\big)^m\\
        \vdots&\vdots&\ddots&\vdots\\
        \sum_{j\in[k]}w_j\big(e^{i\pi\mu^{(1)}}\big)^{m-1}&\sum_{j\in[k]}w_j\big(e^{i\pi\mu^{(1)}}\big)^m&\ldots&\sum_{j\in[k]}w_j\big(e^{i\pi\mu^{(1)}}\big)^{2m-1}\\
    \end{pmatrix}\\
    &=\begin{pmatrix}
        f(0)&f(1)&\ldots&f(m-1)\\
        f(1)&f(2)&\ldots&f(m)\\
        \vdots&\vdots&\ddots&\vdots\\
        f(m-1)&f(m)&\ldots&f(2m-1)
    \end{pmatrix},
\end{flalign*} 
where $f(s)=\sum_{j\in[k]}w_j\big(e^{i\pi\mu^{(j)}}\big)^s$ for $s=0,\ldots,2m-1$ which indeed corresponds to the defined Hankel matrix. The proof is similar for $H_1$ except that $s$ varies from 1 to $2m$ since element-wise multiplication of $H_0$ with the diagonal matrix $D_\mu$ yields a simple scaling by a factor $e^{i\pi\mu^{(j)}}$.

\clearpage
\pagestyle{numberonly}
\printbibliography

\end{document}